\documentclass[11pt,a4paper]{article}

\usepackage[T1]{fontenc}
\usepackage[utf8]{inputenc}
\usepackage{lmodern}
\usepackage{amssymb,amsmath,amsthm,mathtools}
\usepackage[margin=2.5cm]{geometry}
\usepackage{parskip}
\usepackage[hidelinks]{hyperref}
\usepackage{microtype}

% Theorem environments (arXiv / journal style)
\theoremstyle{plain}
\newtheorem{theorem}{Theorem}[section]
\newtheorem{proposition}[theorem]{Proposition}
\newtheorem{lemma}[theorem]{Lemma}
\newtheorem{corollary}[theorem]{Corollary}

\theoremstyle{definition}
\newtheorem{definition}[theorem]{Definition}
\newtheorem{example}[theorem]{Example}
\newtheorem{remark}[theorem]{Remark}

% Notation
\newcommand{\Type}{\mathbf{Type}}
\newcommand{\RTree}{\mathrm{RTree}}
\newcommand{\RPos}{\mathrm{RPos}}
\newcommand{\Fin}{\mathrm{Fin}}

\title{$N$-ary Proof Search Trees\\
  as a Directed Container}
\author{Extracted from a machine-checked development}
\date{}

\begin{document}
\maketitle

\begin{abstract}
We generalise binary proof search trees to finitely branching rose trees
of arbitrary arity and exhibit them as a directed container. Shapes are
rose trees; positions are paths from the root indexed by finite ordinals;
the directed structure is path concatenation. The five directed-container
laws hold by structural recursion on paths. The induced comonad extracts
the root label and duplicates by replacing each label with the full
arbitrarily branching sub-search from that node. A coKleisli morphism is
a proof strategy that may inspect the entire search tree. All stated
results have been machine-checked.
\end{abstract}

\tableofcontents
\bigskip

%======================================================================
\section{Introduction}
%======================================================================

Proof search need not be binary: a tactic may offer any finite number of
alternatives, and dead ends or successes may occur at any arity. Rose
trees---nodes with a finite list of children---are the natural shapes.
We encode branching by a natural number~$n$ together with a family of
children indexed by the finite ordinal $\Fin(n)$, so that structural
recursion works directly on trees and paths without auxiliary
well-founded arguments.

The directed-container structure is identical in spirit to the binary
case: the root is the empty path; subtrees are obtained by descending
along a path; shift is path concatenation. The five laws hold by the
same style of structural recursion.

\paragraph{What this note claims.}
In ordinary mathematical language: finitely branching rose trees with
path positions form a directed container. Binary search trees embed into
this structure (leaves as arity~$0$, binary nodes as arity~$2$), but the
$N$-ary container is not a strict specialisation of the binary one,
since every natural number is allowed as branching factor. The induced
comonad navigates arbitrarily branching search spaces while preserving
dead ends at every depth.

%======================================================================
\section{Shapes and positions}
%======================================================================

\begin{definition}[Rose tree]\label{def:rtree}
A \textbf{finitely branching rose tree} is an element of the inductive
type $\RTree$ generated by a single constructor
\begin{equation}
  \mathrm{node} \;:\;
  (n : \mathbb{N}) \to
  \bigl(\Fin(n) \to \RTree\bigr) \to
  \RTree.
\end{equation}
A leaf (dead end or success) is $\mathrm{node}(0, f)$ where
$f : \Fin(0) \to \RTree$ is vacuous. The only position in a leaf is
$\mathsf{here}$.
\end{definition}

\begin{remark}[Encoding choice]
An encoding via lists of children is more familiar but introduces a
nested inductive occurrence ($\mathrm{List}$ wraps the recursive
argument). The $\Fin(n)$-family encoding avoids that complication and
lets structural recursion work directly on $\RTree$ and $\RPos$.
\end{remark}

\begin{definition}[Position]\label{def:rpos}
A \textbf{position} (path) in a rose tree~$t$ is an element of the
inductive family $\RPos(t)$ generated by
\begin{itemize}
  \item $\mathsf{here}$ --- the root of~$t$;
  \item $\mathsf{child}(i, p)$ --- for $t = \mathrm{node}(n, c)$ and
        $i \in \Fin(n)$, descend into the $i$-th child $c(i)$ and
        continue along path~$p$.
\end{itemize}
For a leaf ($n = 0$), there is no valid $i \in \Fin(0)$, so the only
position is $\mathsf{here}$.
\end{definition}

%======================================================================
\section{Subtree and path concatenation}
%======================================================================

\begin{definition}[Subtree]\label{def:sub}
The \textbf{subtree} operation
$\mathord{\downarrow} : (t : \RTree) \to \RPos(t) \to \RTree$
is defined by
\begin{align}
  t \mathbin{\downarrow} \mathsf{here}
  &\;=\; t, \\
  \mathrm{node}(n, c) \mathbin{\downarrow} \mathsf{child}(i, p)
  &\;=\; c(i) \mathbin{\downarrow} p.
\end{align}
\end{definition}

\begin{definition}[Shift]\label{def:shift}
The \textbf{shift} (path concatenation)
\begin{equation}
  \oplus \;:\;
  (t : \RTree) \to
  (p : \RPos(t)) \to
  \RPos\bigl(t \mathbin{\downarrow} p\bigr) \to
  \RPos(t)
\end{equation}
is defined by
\begin{align}
  \mathsf{here} \oplus q
  &\;=\; q, \\
  \mathsf{child}(i, p) \oplus q
  &\;=\; \mathsf{child}\bigl(i,\; p \oplus q\bigr).
\end{align}
\end{definition}

%======================================================================
\section{The directed-container laws}
%======================================================================

As in the binary case, a directed container is a container with root,
subtree, and shift operations satisfying laws D1--D5. We take
$\mathrm{Shape} = \RTree$, $\mathrm{Pos} = \RPos$,
$o(t) = \mathsf{here}$, and the operations of
Definitions~\ref{def:sub} and~\ref{def:shift}.

\begin{proposition}[D1 and D2, definitional]\label{prop:d1-d2}
For every rose tree~$t$ and every position~$q$ in~$t$,
\begin{equation}
  t \mathbin{\downarrow} \mathsf{here} \;=\; t,
  \qquad
  \mathsf{here} \oplus q \;=\; q.
\end{equation}
\end{proposition}

\begin{proof}
By lean.
\end{proof}

\begin{proposition}[D3: right identity]\label{prop:d3}
For every rose tree~$t$ and every position~$p$ in~$t$,
\begin{equation}
  p \oplus \mathsf{here} \;=\; p.
\end{equation}
\end{proposition}

\begin{proof}
By lean.
\end{proof}

\begin{proposition}[D4: nesting of subtrees]\label{prop:d4}
For every rose tree~$t$, position~$p$ in~$t$, and position~$q$ in the
subtree at~$p$,
\begin{equation}
  t \mathbin{\downarrow} (p \oplus q)
  \;=\;
  \bigl(t \mathbin{\downarrow} p\bigr) \mathbin{\downarrow} q.
\end{equation}
\end{proposition}

\begin{proof}
By lean.
\end{proof}

\begin{proposition}[D5: associativity of shift]\label{prop:d5}
For every rose tree~$t$ and compatible positions $p$, $q$, $r$,
\begin{equation}
  (p \oplus q) \oplus r' \;=\; p \oplus (q \oplus r),
\end{equation}
where $r'$ is $r$ transported along the equality of fibres supplied
by~D4.
\end{proposition}

\begin{proof}
By lean.
\end{proof}

\begin{theorem}[$N$-ary proof search as directed container]
\label{thm:proofsearchn}
The data
\[
  \bigl(
    \RTree,\;
    \RPos,\;
    o = \mathsf{here},\;
    {\mathord{\downarrow}},\;
    {\oplus}
  \bigr)
\]
form a directed container: laws D1--D5 hold.
\end{theorem}

\begin{proof}
By lean.
\end{proof}

%======================================================================
\section{The induced comonad}
%======================================================================

\begin{remark}[Extract and duplicate]
The directed container induces a comonad on the extension. For a
labelling $v : \RPos(t) \to X$,
\begin{align}
  \mathrm{extract}\,\langle t, v \rangle
  &\;=\; v(\mathsf{here}), \\[4pt]
  \mathrm{duplicate}\,\langle t, v \rangle
  &\;=\;
  \Bigl\langle
    t,\;
    p \mapsto
    \bigl\langle
      t \mathbin{\downarrow} p,\;
      q \mapsto v(p \oplus q)
    \bigr\rangle
  \Bigr\rangle.
\end{align}
Thus $\mathrm{extract}$ reads the label at the root, and
$\mathrm{duplicate}$ replaces the label at every node by the full
sub-search from that node, at arbitrary branching factor.
\end{remark}

\begin{remark}[Proof strategies]
A coKleisli morphism $W(X) \to Y$ is a proof strategy that may inspect
the entire (arbitrarily branching) search tree to decide what to
produce.
\end{remark}

%======================================================================
\section{Relation to binary search trees}
%======================================================================

\begin{remark}[Embedding]
Binary search tree shapes embed into rose trees: a leaf maps to
$\mathrm{node}(0, {!})$, and a binary node $\mathsf{node}(\ell, r)$ maps
to $\mathrm{node}\bigl(2,\, [\![\ell', r']\!]\bigr)$, where $\ell'$ and
$r'$ are the recursive images. This is not a strict sub-case of the
binary theory: binary trees force arity $0$ or~$2$ at every node, while
rose trees allow any natural number as branching factor.
\end{remark}

%======================================================================
\section{Discussion}
%======================================================================

\paragraph{What has been established.}
Finitely branching rose trees with path positions form a directed
container. The five laws hold by structural recursion on paths; D1
and~D2 hold definitionally. The induced comonad navigates arbitrarily
branching search spaces while preserving dead ends at every depth.

\paragraph{Machine checking.}
Every definition and theorem statement above corresponds to a declaration
in a machine-checked formalisation. The phrase ``by lean'' indicates that
the proof has been verified by a proof assistant.

%======================================================================
\section*{Acknowledgements}
%======================================================================

The mathematical content follows the theory of directed containers of
Ahman, Chapman, and Uustalu, generalising the binary construction. The
presentation here is a prose rendering of a formal development intended
for human review of the theorem statements.

\begin{thebibliography}{9}

\bibitem{ahman2012}
D.~Ahman, J.~Chapman, and T.~Uustalu.
\newblock When is a container a comonad?
\newblock In \emph{FoSSaCS}, volume~7213 of \emph{LNCS}, pages 74--88.
\newblock Springer, 2012.

\bibitem{abbott2005}
M.~Abbott, T.~Altenkirch, and N.~Ghani.
\newblock Containers: constructing strictly positive types.
\newblock \emph{Theoretical Computer Science}, 342(1):3--27, 2005.

\end{thebibliography}

\end{document}
